\documentclass{article}

\usepackage{amsmath,amssymb}
\usepackage{xurl}
\usepackage[hidelinks]{hyperref}

% Shared commands for notes in docs/paper-gaps/.

\DeclareUnicodeCharacter{2080}{\ifmmode{}_{0}\else\textsubscript{0}\fi}
\DeclareUnicodeCharacter{2081}{\ifmmode{}_{1}\else\textsubscript{1}\fi}
\DeclareUnicodeCharacter{2082}{\ifmmode{}_{2}\else\textsubscript{2}\fi}
\DeclareUnicodeCharacter{2099}{\ifmmode{}_{n}\else\textsubscript{n}\fi}

\newcommand{\Lean}{\textsc{Lean}}
\ExplSyntaxOn
\NewDocumentCommand{\leanid}{m}{
  \ifmmode
    \textsf{\detokenize{#1}}
  \else
    \group_begin:
    \urlstyle{sf}
    \tl_set:Nn \l_tmpa_tl {#1}
    \tl_replace_all:Nnn \l_tmpa_tl {\_} {_}
    \exp_args:NV \path \l_tmpa_tl
    \group_end:
  \fi
}
\ExplSyntaxOff
\providecommand{\tr}{\operatorname{tr}}
\newcommand{\tnleanIssueBase}{https://github.com/LionSR/TNLean/issues/}
\newcommand{\tnleanPrBase}{https://github.com/LionSR/TNLean/pull/}
\newcommand{\tnleanDocBase}{https://sirui-lu.com/TNLean/docs/}
\newcommand{\ghissue}[1]{\href{\tnleanIssueBase#1}{GitHub issue \##1}}
\newcommand{\ghpr}[1]{\href{\tnleanPrBase#1}{PR \##1}}
\newcommand{\doclink}[1]{\href{\tnleanDocBase#1.html}{\path{#1}}}

% Dirac notation and the identity operator, as in blueprint/src/macros/common.tex.
% \providecommand keeps note-local definitions authoritative.
\usepackage{dsfont}
\providecommand{\ket}[1]{|#1\rangle}
\providecommand{\bra}[1]{\langle #1|}
\providecommand{\braket}[2]{\langle #1 | #2 \rangle}
\providecommand{\ketbra}[2]{|#1\rangle\!\langle #2|}
\providecommand{\Id}{\mathds{1}}
\providecommand{\one}{\mathds{1}}
\providecommand{\C}{\mathbb{C}}
\providecommand{\MN}[1]{M_{#1}(\C)}
\providecommand{\MD}{M_D(\C)}

% External referencing for paper sources.  The build helper writes sanitized
% auxiliary files under build/paper-aux/ and excludes duplicated source labels.
\usepackage{xr}

% Import a generated paper auxiliary file with a caller-supplied label prefix.
% The candidate paths support builds from docs/paper-gaps/, the repository root,
% and build/paper-aux/ itself.
\newcommand{\importpaperaux}[2]{%
  \IfFileExists{../../build/paper-aux/#2.aux}{%
    \externaldocument[#1]{../../build/paper-aux/#2}%
  }{%
    \IfFileExists{build/paper-aux/#2.aux}{%
      \externaldocument[#1]{build/paper-aux/#2}%
    }{%
      \IfFileExists{#2.aux}{%
        \externaldocument[#1]{#2}%
      }{}%
    }%
  }%
}

% General external reference macros
\newcommand{\paperref}[2]{\ref{#1-#2}}
\newcommand{\papereqref}[2]{(\ref{#1-#2})}

% CPSV16 source (1606.00608 / MPDO-22-12-17-2.tex) external document and shortcuts
\importpaperaux{cpsv16-}{1606.00608/MPDO-22-12-17-2-tnlean}
\newcommand{\cpsvref}[1]{\paperref{cpsv16}{#1}}
\newcommand{\cpsveqref}[1]{\papereqref{cpsv16}{#1}}

% Verdict marker: \gapnote{<kind>}{<status>}. Kind is the policy
% classification (clarification, local-correction, scope-restriction,
% unfaithful, false-source, open-gap); status is open, wip (elimination
% underway), resolved, or historical. Machine-read by the site index;
% prints a verdict line.
\newcommand{\gapnote}[2]{\par\noindent\textbf{Verdict:} #1 (#2).\par}


\title{The First MPU Source Cut and the Orientation of the Source Gates}
\author{}
\date{2026-09-03}

\begin{document}
\maketitle
\gapnote{local-correction}{resolved}

\section{Source-cut orientation}

Let $\mathcal U^i_{j,\alpha,\beta}$ be a matrix product operator tensor, with
upper and lower physical indices $i,j$ and left and right auxiliary indices
$\alpha,\beta$. Section~III of Ref.~\cite{Cirac2017MPU} introduces two matrix
flattenings. The diagonal cut in the first diagram separates the upper and
right legs from the left and lower legs. The second cut separates the left and
upper legs from the lower and right legs. Hence
\begin{align}
    (M_1)_{(i,\beta),(\alpha,j)}
     & =\mathcal U^i_{j,\alpha,\beta},
    \label{eq:mpu_source_cut_orientation_first_cut} \\
    (M_2)_{(\alpha,i),(j,\beta)}
     & =\mathcal U^i_{j,\alpha,\beta}.
    \label{eq:mpu_source_cut_orientation_second_cut}
\end{align}
The same leg assignment is explicit in equation \texttt{eq:3-leg} of
Ref.~\cite{FrancoRubio2025MPUGauging}. In particular, $X_1$ carries $(i,\beta)$
and $Y_1$ carries $(\alpha,j)$ in the factorization $M_1=X_1Y_1$.

An earlier formal reading used the transpose of
Eq.~\ref{eq:mpu_source_cut_orientation_first_cut}, with row $(\alpha,j)$ and
column $(i,\beta)$. That reading preserved the rank of $M_1$, but it changed
the types of $X_1$ and $Y_1$ and therefore changed every source contraction in
which their open legs mattered. The corrected source cuts are
\leanid{MPOTensor.sourceCutM₁} and \leanid{MPOTensor.sourceCutM₂}.

\section{Weight and adjunction}

The row space of $M_1$ is ordered as the upper physical factor followed by the
right auxiliary factor. The weighted normalization in
Ref.~\cite{Cirac2017MPU} is therefore represented without an additional swap:
\begin{align}
    W_\rho
     & =\operatorname{Id}_d\otimes\rho,
    \label{eq:mpu_source_cut_orientation_weight} \\
    X_1^\dagger W_\rho X_1
     & =\operatorname{Id}_r.
    \label{eq:mpu_source_cut_orientation_weighted_isometry}
\end{align}
These are \leanid{MPOTensor.sourceWeight} and
\leanid{MPOTensor.sourceX₁_weighted_isometry}. Physical adjunction exchanges
the two cuts by conjugate transpose,
\begin{align}
    M_1(\mathcal U^\sharp)
     & =M_2(\mathcal U)^\dagger,
    \label{eq:mpu_source_cut_orientation_adjoint_one} \\
    M_2(\mathcal U^\sharp)
     & =M_1(\mathcal U)^\dagger.
    \label{eq:mpu_source_cut_orientation_adjoint_two}
\end{align}
The corresponding formal declarations are
\leanid{MPOTensor.sourceCutM₁_physicalAdjointTensor} and
\leanid{MPOTensor.sourceCutM₂_physicalAdjointTensor}.

\section{Source gates}

With $M_1=X_1Y_1$ and $M_2=X_2Y_2$, the gates in equations
\texttt{uuvv} of Ref.~\cite{Cirac2017MPU} and \texttt{eq:uv} of
Ref.~\cite{FrancoRubio2025MPUGauging} have entries
\begin{align}
    u_{(l,r),(i_1,i_2)}
     & =\sum_\beta
    (Y_2)_{l,(i_1,\beta)}(Y_1)_{r,(\beta,i_2)},
    \label{eq:mpu_source_cut_orientation_u} \\
    v_{(j_1,j_2),(r,l)}
     & =\sum_\alpha
    (X_1)_{(j_1,\alpha),r}(X_2)_{(\alpha,j_2),l}.
    \label{eq:mpu_source_cut_orientation_v}
\end{align}
They are formalized by \leanid{MPOTensor.SourceFactors.sourceU} and
\leanid{MPOTensor.SourceFactors.sourceV}. The periodic two-site operator then
factors as
\begin{align}
    U^{(2)}_{(i_1,i_2),(j_1,j_2)}
     & =\sum_{l,r}
    v_{(i_1,i_2),(r,l)}u_{(l,r),(j_2,j_1)}.
    \label{eq:mpu_source_cut_orientation_two_site}
\end{align}
This is \leanid{MPOTensor.mpo_two_pair_entry_eq_sourceV_mul_sourceU_swap}.
The open two-site forms use $X_2,u,X_1$ or $v,Y_1,Y_2$, as recorded by
\leanid{MPOTensor.SourceFactors.mul_apply_eq_sum_X₂_mul_sourceU_mul_X₁} and
\leanid{MPOTensor.SourceFactors.mul_apply_eq_sum_sourceV_mul_Y₁_mul_Y₂}.

The contractions $Y_1$--$X_2$ and $X_1$--$Y_2$ remain useful auxiliary
kernels, but they are not the gates $u$ and $v$. Treating either mixed kernel
as a paper gate produces the wrong physical and auxiliary leg assignment.

\section{Affected claims and audit boundary}

The transposed first cut affected source-factor recovery, physical adjunction,
physical reindexing, tensor products of supplied source factors, open and
closed source networks, source-$v$ dressing, and the explicit shift examples.
Rank-only statements could remain numerically true because transposition
preserves rank, while entrywise and open-leg statements required new proofs.
The corrected declarations now use the leg assignments in
Eqs.~\ref{eq:mpu_source_cut_orientation_first_cut}--
\ref{eq:mpu_source_cut_orientation_v}.

Any source-$u$ or source-$v$ result written before this correction must be
revalidated from its current declaration rather than accepted from an older
compiled artifact. In particular, a proof must expose the $Y_2$--$Y_1$
contraction for $u$ or the $X_1$--$X_2$ contraction for $v$, with the product
orders shown above. The Chapter~28 Blueprint links only to declarations that
were rebuilt after the correction.

\section{Remaining source-$u$ comparison}

The source-$u$ Gram contraction has now been expanded into the staggered
four-tensor network and closed as
\begin{align}
    \sum_{l,r}u_{(l,r),q}\overline{u_{(l,r),p}}
     & =\operatorname{tr}\!\left(
    W^{p_2q_2}\lvert\rho^{\mathsf T})(\Phi\rvert W^{p_1q_1}
    \right).
    \label{eq:mpu_source_cut_orientation_transpose_trace}
\end{align}
The coordinate expansion and this closure are
\leanid{MPOTensor.transpose_fixed_pair_trace_eq_staggered_four_u} and
\leanid{MPOTensor.SourceFactors.sourceU_gram_eq_transpose_fixed_pair_trace}.

The retained input word can also be separated into its two one-site letters
and its blocked interior:
\begin{align}
    W^{(K+2)}_{(p_1,p_2,\sigma),(q_1,q_2,\tau)}
     & =\operatorname{tr}\!\left(
    W^{p_2q_2}W_K^{\sigma,\tau}W^{p_1q_1}\right).
    \label{eq:mpu_source_cut_orientation_endpoint_separated}
\end{align}
This is
\leanid{MPOTensor.mpo_doubleLayer_finAddTwoArrowEquiv_symm_eq_endpointSeparated}.
If $E^J=\lvert\rho)(\Phi\rvert$, $\operatorname{tr}(\rho)=1$, and
$K=JD^2$, the normalized input-tail contraction is
\begin{align}
    d^{-K}\sum_{\tau,\eta}
    \overline{U^{(K+2)}_{\eta,(p,\tau)}}
    U^{(K+2)}_{\eta,(q,\tau)}
     & =\operatorname{tr}\!\left(
    W^{p_2q_2}\lvert\rho)(\Phi\rvert W^{p_1q_1}
    \right).
    \label{eq:mpu_source_cut_orientation_direct_trace}
\end{align}
This is
\leanid{MPOTensor.normalized_mpo_input_tail_eq_fixed_pair_trace}.

Consequently the unresolved step is exactly the scalar equality between
Eqs.~\eqref{eq:mpu_source_cut_orientation_transpose_trace} and
\eqref{eq:mpu_source_cut_orientation_direct_trace}. It compares the
transpose-oriented fixed pair of the staggered source-$u$ Gram network with
the direct fixed pair of the normalized input tail. No symmetry of $\rho$
is among the stated assumptions. Therefore the complete source-$u$
contraction, the source-$u$ isometry, and the subsequent consequences do not
follow from the established identities alone.

\section{Verdict}

The former definition of $M_1$ was the transpose of the source diagram. The
correction changes no source rank, but it changes the typed source factors and
all order-sensitive contractions. The first cut, weight, adjunction identities,
paper gates, two-site factorizations, tensor products, and shift examples now
use one consistent leg assignment. The orientation correction is resolved.
The remaining source-$u$ question is the scalar comparison between the
transpose-oriented and direct fixed-pair traces above.

\bibliographystyle{alpha}
\bibliography{references}

\end{document}
